Smart contracts accept zero-knowledge (ZK) proofs as permission to release
funds, mint credentials and settle state. A proof certifies only its public
statement, yet the contract's action may depend on values the statement never
fixed, such as a recipient taken from calldata. Anyone who observes a valid
proof can then redirect or repeat the effect. We call these \emph{binding
flaws}. Two of them at the proof-to-contract boundary caused about \$7.3M in
losses in 2025--2026. We define binding as a relation between a
\emph{statement ledger} and an \emph{effect ledger} under the application's
intended policy, and classify violations into five patterns that mirror
classic signature misuse. In 8,578 public audit reports we find 87 core
binding flaws across 43 project families; 78\% of the rated ones are Critical
or High, and one in six ZK application audits reports at least one. Static
analyzers, circuit analyzers and generic fuzzing find none of the audited
flaws they can analyze. We present \tool{}, an LLM agent that must fill both
ledgers under coverage checks, and use its ledgers to ask why such an agent
misses flaws. The misses we could recover were not unseen: the agent reported
them on the wrong side of the proof boundary, or recorded them and never
reported them. A tool-less
post-stage that re-adjudicates discarded suspicions under closed
admissibility rules raises strict recall on 28 development flaws from 12 to 18 over two runs
(10.0 to 14.0 per run; paired $p=0.031$) for \$0.15 per scan, with no strict finding on 17 patched
versions, while four knowledge- and tool-based extensions fail a
specificity rule fixed in advance. \DEMO{On 51 held-out flaws the frozen tool gains 10 and
loses 2 against its predecessor ($p=0.04$), and on 25 audited applications
it found three previously unknown binding flaws, confirmed by executed attack
tests and reported to the maintainers.}
